\documentclass{beamer}

\usetheme{Madrid}

\usepackage{amsmath,amssymb,amsfonts,bm}
\usepackage{physics}
\usepackage{mathtools}
\usepackage{graphicx}

\title{The Calogero--Sutherland Hamiltonian}
\author{Morten Hjorth-Jensen}
\date{Spring 2025}

\begin{document}

%------------------------------------------------
\begin{frame}
\titlepage
\end{frame}

%------------------------------------------------
\begin{frame}{Introduction}

The \textbf{Calogero--Sutherland model} is a celebrated exactly solvable
quantum many-body system describing particles interacting through
long-range inverse-square interactions.

\vspace{0.3cm}

It plays important roles in

\begin{itemize}
    \item Integrable quantum systems
    \item Fractional statistics
    \item Random matrix theory
    \item Quantum Hall physics
    \item Conformal field theory
    \item Spin chains and long-range magnets
\end{itemize}

\vspace{0.3cm}

The model exists in several related forms:
\begin{itemize}
    \item Calogero model (particles on a line)
    \item Sutherland model (particles on a circle)
\end{itemize}

\end{frame}

%------------------------------------------------
\begin{frame}{The Calogero Hamiltonian}

The Calogero model describes $N$ particles on a line with harmonic confinement.

\vspace{0.3cm}

The Hamiltonian is

\[
H =
-\frac{\hbar^2}{2m}
\sum_{i=1}^{N}
\frac{\partial^2}{\partial x_i^2}
+
\frac{m\omega^2}{2}
\sum_{i=1}^{N}x_i^2
+
\sum_{i<j}
\frac{g}{(x_i-x_j)^2}.
\]

\vspace{0.3cm}

The three terms correspond to

\begin{enumerate}
    \item Kinetic energy
    \item Harmonic trapping potential
    \item Inverse-square interaction
\end{enumerate}

\end{frame}

%------------------------------------------------
\begin{frame}{Coupling Constant Parametrization}

A common parametrization is

\[
g =
\frac{\hbar^2}{m}
\lambda(\lambda-1),
\]

which gives

\[
H =
-\frac{\hbar^2}{2m}
\sum_i \partial_i^2
+
\frac{m\omega^2}{2}
\sum_i x_i^2
+
\frac{\hbar^2}{m}
\sum_{i<j}
\frac{\lambda(\lambda-1)}
{(x_i-x_j)^2}.
\]

\vspace{0.4cm}

The parameter $\lambda$ controls:
\begin{itemize}
    \item Interaction strength
    \item Correlation effects
    \item Generalized statistics
\end{itemize}

\end{frame}

%------------------------------------------------
\begin{frame}{The Sutherland Model}

The Sutherland model places particles on a ring of circumference $L$.

\vspace{0.3cm}

The Hamiltonian becomes

\[
H =
-\frac{\hbar^2}{2m}
\sum_{i=1}^{N}
\frac{\partial^2}{\partial x_i^2}
+
\frac{\hbar^2}{m}
\left(\frac{\pi}{L}\right)^2
\sum_{i<j}
\frac{\lambda(\lambda-1)}
{\sin^2\left(
\frac{\pi(x_i-x_j)}{L}
\right)}.
\]

\vspace{0.3cm}

Important features:
\begin{itemize}
    \item Periodic geometry
    \item Translational invariance
    \item Exact solvability
    \item Long-range interactions
\end{itemize}

\end{frame}

%------------------------------------------------
\begin{frame}{Ground-State Wavefunction}

The Calogero ground state is exactly known:

\[
\Psi_0(x_1,\dots,x_N)
=
\prod_{i<j}
|x_i-x_j|^\lambda
\exp\left(
-\frac{m\omega}{2\hbar}
\sum_i x_i^2
\right).
\]

\vspace{0.4cm}

The factor

\[
\prod_{i<j}|x_i-x_j|^\lambda
\]

is called a \textbf{Jastrow factor}.

\vspace{0.3cm}

It produces strong correlations between particles.

\end{frame}

%------------------------------------------------
\begin{frame}{Jastrow Correlations}

The Jastrow structure appears throughout many-body physics:

\[
\Psi \sim
\prod_{i<j}
f(x_i-x_j).
\]

\vspace{0.3cm}

Connections include:

\begin{itemize}
    \item Laughlin quantum Hall states
    \item Variational Monte Carlo
    \item Neural quantum states
    \item Tensor-network-inspired wavefunctions
\end{itemize}

\vspace{0.3cm}

The inverse-square repulsion suppresses particle overlap.

\end{frame}

%------------------------------------------------
\begin{frame}{Integrability}

The Calogero--Sutherland model is integrable.

\vspace{0.3cm}

Important consequences:

\begin{itemize}
    \item Exact spectrum
    \item Infinite set of conserved quantities
    \item Exact eigenfunctions
    \item Solvable dynamics
\end{itemize}

\vspace{0.4cm}

The classical version admits a Lax-pair formulation:

\[
\dot{L} = [L,M].
\]

\vspace{0.3cm}

This structure is central in integrable systems theory.

\end{frame}

%------------------------------------------------
\begin{frame}{Connection to Random Matrix Theory}

The ground-state probability density is

\[
|\Psi_0|^2
\propto
\prod_{i<j}
|x_i-x_j|^{2\lambda}.
\]

\vspace{0.3cm}

This is precisely the structure appearing in Dyson ensembles
of random matrix theory.

\vspace{0.3cm}

Connections:
\begin{itemize}
    \item Gaussian Orthogonal Ensemble (GOE)
    \item Gaussian Unitary Ensemble (GUE)
    \item Gaussian Symplectic Ensemble (GSE)
\end{itemize}

\vspace{0.3cm}

The particles behave similarly to repelling eigenvalues.

\end{frame}

%------------------------------------------------
\begin{frame}{Fractional Statistics}

The parameter $\lambda$ interpolates between different statistics.

\vspace{0.4cm}

Examples:

\[
\lambda = 0
\quad \Rightarrow \quad
\text{Bosonic limit}
\]

\vspace{0.3cm}

\[
\lambda = 1
\quad \Rightarrow \quad
\text{Fermionic-like behavior}
\]

\vspace{0.3cm}

General $\lambda$ corresponds to generalized exclusion statistics.

\vspace{0.3cm}

This connects the model to:
\begin{itemize}
    \item Anyons
    \item Fractional quantum Hall systems
    \item Haldane statistics
\end{itemize}

\end{frame}

%------------------------------------------------
\begin{frame}{Second-Quantized Viewpoint}

The interaction potential is

\[
V(x_i-x_j)
=
\frac{g}{(x_i-x_j)^2}.
\]

\vspace{0.3cm}

This is highly nonlocal compared to:

\begin{itemize}
    \item Contact interactions
    \[
    g\,\delta(x_i-x_j)
    \]
    
    \item Nearest-neighbor lattice interactions
\end{itemize}

\vspace{0.3cm}

The long-range structure produces highly correlated many-body states.

\end{frame}

%------------------------------------------------
\begin{frame}{Connections to Quantum Hall Physics}

The Jastrow structure resembles Laughlin states:

\[
\Psi_{\mathrm{Laughlin}}
\sim
\prod_{i<j}
(z_i-z_j)^m
e^{-\sum_i |z_i|^2/4}.
\]

\vspace{0.3cm}

Both systems exhibit:

\begin{itemize}
    \item Strong correlations
    \item Fractional statistics
    \item Collective excitations
    \item Emergent many-body structure
\end{itemize}

\end{frame}

%------------------------------------------------
\begin{frame}{Modern Applications}

The Calogero--Sutherland model appears in modern research on:

\begin{itemize}
    \item Neural quantum states
    \item Variational quantum eigensolvers (VQE)
    \item Quantum simulation
    \item Long-range spin chains
    \item Random matrices
    \item Quantum information
    \item Diffusion-based quantum state generation
\end{itemize}

\vspace{0.3cm}

It serves as an important benchmark for correlated quantum systems.

\end{frame}

%------------------------------------------------
\begin{frame}{Summary}

The Calogero--Sutherland model is a fundamental exactly solvable
many-body quantum system.

\vspace{0.4cm}

Key features:

\begin{itemize}
    \item Inverse-square interactions
    \item Exact solvability
    \item Jastrow-correlated wavefunctions
    \item Fractional statistics
    \item Deep links to random matrices and quantum Hall physics
\end{itemize}

\vspace{0.4cm}

It continues to play an important role in:
\begin{itemize}
    \item Condensed matter physics
    \item Mathematical physics
    \item Quantum computing
    \item Machine learning for quantum systems
\end{itemize}

\end{frame}

%------------------------------------------------
\end{document}
